      Solutions for the first day problems at the IMC 2000

     Problem 1.
     Is it true that if f : [0, 1] → [0, 1] is
     a) monotone increasing
     b) monotone decreasing
     then there exists an x ∈ [0, 1] for which f (x) = x?
     Solution.
     a) Yes.
     Proof: Let A = {x ∈ [0, 1] : f (x) > x}. If f (0) = 0 we are done, if not then A is
non-empty (0 is in A) bounded, so it has supremum, say a. Let b = f (a).
     I. case: a < b. Then, using that f is monotone and a was the sup, we get b = f (a) ≤
f ((a + b)/2) ≤ (a + b)/2, which contradicts a < b.
     II. case: a > b. Then we get b = f (a) ≥ f ((a + b)/2) > (a + b)/2 contradiction.
Therefore we must have a = b.
     b) No. Let, for example,

                                 f (x) = 1 − x/2 if     x ≤ 1/2
and
                                f (x) = 1/2 − x/2 if     x > 1/2
      This is clearly a good counter-example.

     Problem 2.
     Let p(x) = x5 + x and q(x) = x5 + x2 . Find all pairs (w, z) of complex numbers with
w 6= z for which p(w) = p(z) and q(w) = q(z).
     Short solution. Let

                              p(x) − p(y)
                 P (x, y) =               = x4 + x3 y + x2 y 2 + xy 3 + y 4 + 1
                                 x−y

and
                           q(x) − q(y)
               Q(x, y) =               = x4 + x3 y + x2 y 2 + xy 3 + y 4 + x + y.
                              x−y
We need those pairs (w, z) which satisfy P (w, z) = Q(w, z) = 0.
    From P − Q = 0 we have w + z = 1. Let c = wz. After a short calculation we obtain
c2 − 3c + 2 = 0, which has the solutions c = 1 and c = 2. From the system w + z = 1,
wz = c we obtain the following pairs:
                         √      √ !                        √      √ !
                     1 ± 3i 1 ∓ 3i                     1 ± 7i 1 ∓ 7i
                           ,                  and            ,        .
                        2      2                          2      2




                                                1
    Problem 3.
    A and B are square complex matrices of the same size and

                                         rank(AB − BA) = 1 .

Show that (AB − BA)2 = 0.
    Let C = AB − BA. Since rank C = 1, at most one eigenvalue of C is different from 0.
Also tr C = 0, so all the eigevalues are zero. In the Jordan canonical form there can only
be one 2 × 2 cage and thus C 2 = 0.

    Problem 4.
    a) Show that if (xi ) is a decreasing sequence of positive numbers then

                                             n
                                                         !1/2      n
                                             X                    X   xi
                                                   x2i          ≤     √ .
                                             i=1                  i=1
                                                                       i

   b) Show that there is a constant C so that if (xi ) is a decreasing sequence of positive
numbers then                                 !1/2
                            ∞         ∞                  ∞
                           X    1    X
                                           2
                                                        X
                               √         x        ≤C        xi .
                           m=1
                                 m i=m i                i=1

Solution.
   a)
                n         n        n       i       n          n
               X   xi 2 X xi xj   X   xi X xi     X   xi xi  X
             (     √ ) =     √√ ≥     √       √ ≥     √ i√ =     x2i
               i=1
                    i    i,j
                              i j i=1
                                        i j=1
                                               j  i=1
                                                       i  i  i=1

    b)
                     ∞       ∞          ∞       ∞
                     X    1 X 2 1/2    X    1 X       xi
                         √ (    xi ) ≤     √      √
                     m=1
                          m i=m        m=1
                                             m i=m i − m + 1

by a)
                                     ∞             i
                                     X             X       1
                                 =            xi       √ √
                                     i=1           m=1
                                                        m i−m+1

You can get a sharp bound on

                                              i
                                              X       1
                                     sup          √ √
                                         i    m=1
                                                   m i−m+1

by checking that it is at most
                                 Z i+1
                                                   1
                                               √ √      dx = π
                                     0          x i+1−x

                                                           2
Alternatively you can observe that

                     i                i/2
                     X        1       X       1
                         √  √      =2     √ √      ≤
                     m=1
                           m i+1−m    m=1
                                           m i+1−m

                                     i/2
                             1 X 1         1   p
                         ≤ 2p       √ ≤ 2 p .2 i/2 = 4
                             i/2 m=1 m     i/2


     Problem 5.
     Let R be a ring of characteristic zero (not necessarily commutative). Let e, f and g
be idempotent elements of R satisfying e + f + g = 0. Show that e = f = g = 0.
     (R is of characteristic zero means that, if a ∈ R and n is a positive integer, then
na 6= 0 unless a = 0. An idempotent x is an element satisfying x = x2 .)
     Solution. Suppose that e + f + g = 0 for given idempotents e, f, g ∈ R. Then

               g = g 2 = (−(e + f ))2 = e + (ef + f e) + f = (ef + f e) − g,

i.e. ef+fe=2g, whence the additive commutator

            [e, f ] = ef − f e = [e, ef + f e] = 2[e, g] = 2[e, −e − f ] = −2[e, f ],

i.e. ef = f e (since R has zero characteristic). Thus ef + f e = 2g becomes ef = g, so that
e + f + ef = 0. On multiplying by e, this yields e + 2ef = 0, and similarly f + 2ef = 0,
so that f = −2ef = e, hence e = f = g by symmetry. Hence, finaly, 3e = e + f + g = 0,
i.e. e = f = g = 0.
      For part (i) just omit some of this.

    Problem 6.
    Let f : R → (0, ∞) be an increasing differentiable function for which lim f (x) = ∞
                                                                                  x→∞
and f 0 is bounded.
                 Rx
    Let F (x) = f . Define the sequence (an ) inductively by
                0

                                                            1
                               a0 = 1,     an+1 = an +           ,
                                                         f (an )

and the sequence (bn ) simply by bn = F −1 (n). Prove that lim (an − bn ) = 0.
                                                           n→∞
    Solution. From the conditions it is obvious that F is increasing and lim bn = ∞.
                                                                                  n→∞
     By Lagrange’s theorem and the recursion in (1), for all k ≥ 0 integers there exists a
real number ξ ∈ (ak , ak+1 ) such that

                                                                      f (ξ)
                      F (ak+1 ) − F (ak ) = f (ξ)(ak+1 − ak ) =              .          (2)
                                                                     f (ak )

                                               3
By the monotonity, f (ak ) ≤ f (ξ) ≤ f (ak+1 ), thus
                                                  f (ak+1 )     f (ak+1 ) − f (ak )
                  1 ≤ F (ak+1 ) − F (ak ) ≤                 =1+                     .                     (3)
                                                   f (ak )            f (ak )
Summing (3) for k = 0, . . . , n − 1 and substituting F (bn ) = n, we have
                                                                            n−1
                                                                            X     f (ak+1 ) − f (ak )
          F (bn ) < n + F (a0 ) ≤ F (an ) ≤ F (bn ) + F (a0 ) +                                       .   (4)
                                                                                        f (ak )
                                                                            k=0

From the first two inequalities we already have an > bn and lim an = ∞.
                                                                               n→∞
     Let ε be an arbitrary positive number. Choose an integer Kε such that f (aKε ) > 2ε .
If n is sufficiently large, then
                                              n−1
                                              X     f (ak+1 ) − f (ak )
                                 F (a0 ) +                              =
                                                          f (ak )
                                              k=0

                              K                                      n−1
                                                             !
                                ε −1
                              X        f (ak+1 ) − f (ak )           X      f (ak+1 ) − f (ak )
              =   F (a0 ) +                                      +                              <         (5)
                                             f (ak )                              f (ak )
                              k=0                                    k=Kε
                                                    n−1
                                      1     X                    
                        < Oε (1) +            f (ak+1 ) − f (ak ) <
                                   f (aKε )
                                                    k=Kε
                                     ε                 
                          < Oε (1) + f (an ) − f (aKε ) < εf (an ).
                                     2
Inequalities (4) and (5) together say that for any positive ε, if n is sufficiently large,
                                        F (an ) − F (bn ) < εf (an ).
Again, by Lagrange’s theorem, there is a real number ζ ∈ (bn , an ) such that
                     F (an ) − F (bn ) = f (ζ)(an − bn ) > f (bn )(an − bn ),                             (6)
thus
                                        f (bn )(an − bn ) < εf (an ).                                     (7)
Let B be an upper bound for f 0 . Apply f (an ) < f (bn ) + B(an − bn ) in (7):
                                                                 
                      f (bn )(an − bn ) < ε f (bn ) + B(an − bn ) ,
                                         
                             f (bn ) − εB (an − bn ) < εf (bn ).                                          (8)
Due to lim f (bn ) = ∞, the first factor is positive, and we have
        n→∞

                                                        f (bn )
                                    an − b n < ε                 < 2ε                                     (9)
                                                    f (bn ) − εB
for sufficiently large n.
     Thus, for arbitrary positive ε we proved that 0 < an − bn < 2ε if n is sufficiently large.

                                                      4
